\documentclass[aspectratio=169]{beamer}

\usetheme{Madrid}
\usecolortheme{default}

\usepackage{amsmath}
\usepackage{booktabs}
\usepackage{tikz}
\usepackage{graphicx}
\usepackage{hyperref}

\title{Modern Foundation Models}
\subtitle{Lecture 3: Large Language Models, MoE, Decoding and Prompting}
\author{Dror Lederman}
\institute{Holon Institute of Technology}
\date{\today}

\begin{document}

%------------------------------------------------
\begin{frame}
\titlepage
\end{frame}

%------------------------------------------------
\begin{frame}{Lecture Roadmap}
\begin{enumerate}
    \item Recap: Transformer Model Families
    \item What is a Large Language Model?
    \item Mixture-of-Experts LLMs
    \item Response Generation
    \item Decoding Strategies
    \item Guided Decoding
    \item Prompting Strategies
    \item Inference Optimizations
\end{enumerate}
\end{frame}

%------------------------------------------------
\section{Recap}

\begin{frame}{Recap: Transformer-Based Model Families}

\begin{table}
\centering
\begin{tabular}{lll}
\toprule
Family & Main Use & Examples \\
\midrule
Encoder-Decoder & Text-to-text & T5, mT5, ByT5 \\
Encoder-Only & Understanding, classification & BERT, RoBERTa \\
Decoder-Only & Text generation & GPT, LLaMA, Mistral \\
\bottomrule
\end{tabular}
\end{table}

\vspace{0.4cm}

Modern LLMs are usually \textbf{decoder-only Transformer-based models}.

\end{frame}

%------------------------------------------------
\section{LLM Overview}

\begin{frame}{Terminology}

\[
\text{LLM} = \text{Large Language Model}
\]

\vspace{0.5cm}

A language model assigns probabilities to sequences of tokens:

\[
P(w_1, w_2, \ldots, w_T)
\]

\vspace{0.4cm}

Using the chain rule:

\[
P(w_1,\ldots,w_T)
=
\prod_{t=1}^{T}
P(w_t \mid w_1,\ldots,w_{t-1})
\]

\end{frame}

%------------------------------------------------
\begin{frame}{What Makes a Language Model ``Large''?}

\begin{itemize}
    \item \textbf{Model size:} billions or trillions of parameters.
    \item \textbf{Training data:} hundreds of billions or trillions of tokens.
    \item \textbf{Compute:} large-scale GPU/TPU clusters.
    \item \textbf{Capabilities:} generation, reasoning, coding, summarization, translation.
\end{itemize}

\vspace{0.4cm}

Large scale changes both performance and behavior.

\end{frame}

%------------------------------------------------
\begin{frame}{Decoder-Only LLMs}

Decoder-only LLMs predict the next token:

\[
P(w_{t+1} \mid w_1,\ldots,w_t)
\]

\vspace{0.4cm}

Examples:

\begin{itemize}
    \item GPT series
    \item LLaMA
    \item Gemma
    \item Mistral
    \item Qwen
    \item DeepSeek
\end{itemize}

\end{frame}

%------------------------------------------------
\section{Mixture of Experts}

\begin{frame}{Motivation for Mixture-of-Experts}

Large dense models use all parameters for every token.

\vspace{0.4cm}

Question:

\begin{quote}
Do we really need to activate the entire model for every input?
\end{quote}

\vspace{0.4cm}

Mixture-of-Experts introduces conditional computation:

\[
\text{Use only a subset of the model for each token.}
\]

\end{frame}

%------------------------------------------------
\begin{frame}{Mixture-of-Experts: Basic Idea}

An MoE layer contains multiple experts:

\[
E_1, E_2, \ldots, E_n
\]

A gating network selects which experts to use:

\[
G(x)
\]

The output is:

\[
\hat{y}
=
\sum_{i=1}^{n}
G_i(x)E_i(x)
\]

\end{frame}

%------------------------------------------------
\begin{frame}{Dense MoE vs Sparse MoE}

\textbf{Dense MoE}

\[
\hat{y}
=
\sum_{i=1}^{n}
G_i(x)E_i(x)
\]

All experts contribute.

\vspace{0.5cm}

\textbf{Sparse MoE}

\[
\hat{y}
=
\sum_{i \in \mathcal{I}_k}
G_i(x)E_i(x)
\]

Only top-$k$ experts contribute.

\end{frame}

%------------------------------------------------
\begin{frame}{MoE in Transformer Models}

In modern Transformer-based MoE models:

\begin{itemize}
    \item The dense feed-forward network is replaced by multiple experts.
    \item A router selects experts for each token.
    \item Routing is usually token-level.
\end{itemize}

\vspace{0.4cm}

\[
\text{Transformer Block}
=
\text{Attention}
+
\text{MoE Feed-Forward Layer}
\]

\end{frame}

%------------------------------------------------
\begin{frame}{Why MoE?}

Advantages:

\begin{itemize}
    \item More parameters without proportional increase in compute.
    \item Conditional specialization of experts.
    \item Better scaling efficiency.
\end{itemize}

\vspace{0.4cm}

Tradeoff:

\begin{itemize}
    \item More complex training.
    \item Routing instability.
    \item Load balancing challenges.
\end{itemize}

\end{frame}

%------------------------------------------------
\begin{frame}{Training Challenge: Routing Collapse}

A common problem:

\begin{quote}
The same expert gets selected most of the time.
\end{quote}

\vspace{0.4cm}

This is called:

\[
\text{routing collapse}
\]

\vspace{0.4cm}

It reduces expert diversity and wastes capacity.

\end{frame}

%------------------------------------------------
\begin{frame}{Auxiliary Load-Balancing Loss}

To encourage balanced expert usage:

\[
\mathcal{L}_{aux}
=
\alpha \cdot N
\sum_{i=1}^{N}
f_i P_i
\]

where:

\begin{itemize}
    \item $f_i$ = fraction of tokens routed to expert $i$
    \item $P_i$ = average routing probability for expert $i$
    \item $N$ = number of experts
\end{itemize}

\end{frame}

%------------------------------------------------
\section{Response Generation}

\begin{frame}{Next Token Prediction}

LLMs generate text one token at a time.

\vspace{0.4cm}

Example:

\[
\texttt{[BOS]}
\rightarrow
\texttt{A}
\rightarrow
\texttt{teddy}
\rightarrow
\texttt{bear}
\rightarrow
\texttt{is}
\rightarrow
?
\]

\vspace{0.4cm}

At every step, the model predicts a distribution over the vocabulary.

\end{frame}

%------------------------------------------------
\begin{frame}{Output Probabilities}

The model produces logits:

\[
z \in \mathbb{R}^{|V|}
\]

where $|V|$ is the vocabulary size.

\vspace{0.4cm}

Softmax converts logits into probabilities:

\[
P(w_i \mid C)
=
\frac{\exp(z_i)}
{\sum_{j=1}^{|V|}\exp(z_j)}
\]

\end{frame}

%------------------------------------------------
\begin{frame}{Greedy Decoding}

Greedy decoding chooses the most likely token:

\[
\hat{w}_{t+1}
=
\arg\max_{w}
P(w \mid C)
\]

\vspace{0.4cm}

Advantages:

\begin{itemize}
    \item Simple
    \item Fast
    \item Deterministic
\end{itemize}

Limitations:

\begin{itemize}
    \item Can be repetitive
    \item Often lacks diversity
    \item Locally optimal, not globally optimal
\end{itemize}

\end{frame}

%------------------------------------------------
\begin{frame}{Beam Search}

Beam search keeps the top $k$ most likely partial sequences.

\vspace{0.4cm}

At each step:

\begin{enumerate}
    \item Expand each candidate sequence.
    \item Score all candidates.
    \item Keep only the best $k$ paths.
\end{enumerate}

\vspace{0.4cm}

Beam search approximates:

\[
\arg\max_{w_{1:T}}
P(w_{1:T})
\]

\end{frame}

%------------------------------------------------
\begin{frame}{Beam Search: Pros and Cons}

\begin{table}
\centering
\begin{tabular}{ll}
\toprule
Pros & Cons \\
\midrule
Better than greedy for structured tasks & More expensive \\
Useful for translation/summarization & Often less creative \\
Keeps multiple hypotheses & Can produce generic text \\
\bottomrule
\end{tabular}
\end{table}

\end{frame}

%------------------------------------------------
\begin{frame}{Sampling}

Instead of taking the maximum, sample from the distribution:

\[
\hat{w}_{t+1}
\sim
P(w_{t+1} \mid C)
\]

\vspace{0.4cm}

Sampling produces more diverse outputs.

\vspace{0.4cm}

But uncontrolled sampling may produce low-quality or incoherent text.

\end{frame}

%------------------------------------------------
\begin{frame}{Top-$k$ Sampling}

Top-$k$ sampling restricts sampling to the $k$ most likely tokens.

\vspace{0.4cm}

\[
\mathcal{V}_k
=
\text{top-}k
\left(
P(w \mid C)
\right)
\]

\[
\hat{w}
\sim
P(w \mid C), \quad w \in \mathcal{V}_k
\]

\vspace{0.4cm}

Typical values: $k=40$, $k=50$, $k=100$.

\end{frame}

%------------------------------------------------
\begin{frame}{Top-$p$ Sampling}

Top-$p$ sampling, also called nucleus sampling, chooses the smallest set of tokens whose cumulative probability exceeds $p$.

\[
\sum_{w \in \mathcal{V}_p}
P(w \mid C)
\geq p
\]

\vspace{0.4cm}

Then sample only from $\mathcal{V}_p$.

\vspace{0.4cm}

Typical values:

\[
p = 0.9,\quad p=0.95
\]

\end{frame}

%------------------------------------------------
\begin{frame}{Temperature}

Temperature modifies the sharpness of the probability distribution:

\[
P_T(w_i \mid C)
=
\frac{\exp(z_i/T)}
{\sum_j \exp(z_j/T)}
\]

\vspace{0.4cm}

Interpretation:

\begin{itemize}
    \item $T < 1$: sharper, more deterministic.
    \item $T = 1$: original distribution.
    \item $T > 1$: flatter, more random.
\end{itemize}

\end{frame}

%------------------------------------------------
\begin{frame}{Effect of Temperature}

\begin{table}
\centering
\begin{tabular}{lll}
\toprule
Temperature & Distribution & Output Behavior \\
\midrule
Low & Peaked & Focused, deterministic \\
Medium & Balanced & Natural variation \\
High & Flat & Creative, risky \\
\bottomrule
\end{tabular}
\end{table}

\vspace{0.4cm}

In production systems, temperature is often tuned by task.

\end{frame}

%------------------------------------------------
\section{Guided Decoding}

\begin{frame}{Motivation for Guided Decoding}

Sometimes we need the model to generate output in a strict format.

\vspace{0.4cm}

Example:

\begin{verbatim}
{
  "first_name": "teddy",
  "last_name": "bear",
  "age": 33,
  "hobby": "reading"
}
\end{verbatim}

\vspace{0.4cm}

Prompting alone may not guarantee valid JSON.

\end{frame}

%------------------------------------------------
\begin{frame}{Guided Decoding}

Core idea:

\begin{quote}
Only allow tokens that keep the output valid.
\end{quote}

\vspace{0.4cm}

At each decoding step:

\begin{enumerate}
    \item Compute next-token probabilities.
    \item Mask invalid tokens.
    \item Renormalize probabilities.
    \item Decode from the constrained distribution.
\end{enumerate}

\end{frame}

%------------------------------------------------
\begin{frame}{Token Masking}

Let $\mathcal{A}_t$ be the set of valid next tokens at time $t$.

\[
P'(w \mid C)
=
\begin{cases}
\frac{P(w \mid C)}
{\sum_{u \in \mathcal{A}_t}P(u \mid C)}
&
w \in \mathcal{A}_t
\\
0
&
w \notin \mathcal{A}_t
\end{cases}
\]

\vspace{0.4cm}

This forces syntactic validity.

\end{frame}

%------------------------------------------------
\section{Prompting}

\begin{frame}{Prompting Terminology}

The prompt is the context given to the model.

\vspace{0.4cm}

Important terms:

\begin{itemize}
    \item Context length
    \item Context window
    \item System prompt
    \item User prompt
    \item Few-shot examples
\end{itemize}

\vspace{0.4cm}

The model conditions generation on all tokens in the context window.

\end{frame}

%------------------------------------------------
\begin{frame}{Zero-Shot Prompting}

The model receives only the task instruction.

\vspace{0.4cm}

Example:

\begin{quote}
Classify the sentiment of the following review as positive or negative.
\end{quote}

\vspace{0.4cm}

No examples are provided.

\end{frame}

%------------------------------------------------
\begin{frame}{Few-Shot Prompting}

Few-shot prompting provides examples in the prompt.

\vspace{0.4cm}

\begin{quote}
Review: ``Great product.'' Sentiment: Positive

Review: ``Very disappointing.'' Sentiment: Negative

Review: ``I would buy it again.'' Sentiment:
\end{quote}

\vspace{0.4cm}

The model learns the task pattern from context.

\end{frame}

%------------------------------------------------
\begin{frame}{Chain-of-Thought Prompting}

Chain-of-thought prompting asks the model to reason step by step.

\vspace{0.4cm}

Useful for:

\begin{itemize}
    \item Arithmetic reasoning
    \item Logical reasoning
    \item Multi-step planning
    \item Complex question answering
\end{itemize}

\vspace{0.4cm}

Important caveat: reasoning traces may be unreliable or post-hoc.

\end{frame}

%------------------------------------------------
\begin{frame}{Prompting Is Programming by Context}

Prompting can specify:

\begin{itemize}
    \item Role
    \item Task
    \item Constraints
    \item Output format
    \item Examples
    \item Evaluation criteria
\end{itemize}

\vspace{0.4cm}

Prompt design affects both performance and reliability.

\end{frame}

%------------------------------------------------
\section{Inference Optimizations}

\begin{frame}{Why Inference Optimization Matters}

LLM inference is expensive because:

\begin{itemize}
    \item Models are large.
    \item Generation is sequential.
    \item Each new token depends on previous tokens.
    \item Long contexts increase attention cost.
\end{itemize}

\vspace{0.4cm}

Optimization is crucial for latency, throughput and cost.

\end{frame}

%------------------------------------------------
\begin{frame}{KV Cache}

During autoregressive decoding, keys and values from previous tokens can be reused.

\vspace{0.4cm}

Without cache:

\[
\text{Recompute } K,V \text{ for all previous tokens at every step.}
\]

With cache:

\[
\text{Store previous } K,V \text{ and compute only the new token.}
\]

\vspace{0.4cm}

KV caching greatly improves decoding efficiency.

\end{frame}

%------------------------------------------------
\begin{frame}{Batching}

Batching processes multiple requests together.

\vspace{0.4cm}

Benefits:

\begin{itemize}
    \item Better GPU utilization.
    \item Higher throughput.
    \item Lower cost per generated token.
\end{itemize}

\vspace{0.4cm}

Challenge:

\begin{itemize}
    \item Requests have different lengths and finish at different times.
\end{itemize}

\end{frame}

%------------------------------------------------
\begin{frame}{Quantization}

Quantization reduces numerical precision.

\vspace{0.4cm}

Examples:

\begin{itemize}
    \item FP16
    \item BF16
    \item INT8
    \item INT4
\end{itemize}

\vspace{0.4cm}

Tradeoff:

\[
\text{lower memory and faster inference}
\quad
\leftrightarrow
\quad
\text{possible quality degradation}
\]

\end{frame}

%------------------------------------------------
\begin{frame}{Speculative Decoding}

Speculative decoding uses a small draft model to propose tokens.

\vspace{0.4cm}

Process:

\begin{enumerate}
    \item Small model proposes several tokens.
    \item Large model verifies them in parallel.
    \item Accepted tokens are emitted.
\end{enumerate}

\vspace{0.4cm}

Goal: reduce latency while preserving output distribution.

\end{frame}

%------------------------------------------------
\begin{frame}{Key Takeaways}

\begin{itemize}
    \item Modern LLMs are usually decoder-only Transformers.
    \item MoE increases parameter count using conditional computation.
    \item LLMs generate text by next-token prediction.
    \item Decoding strategy strongly affects output quality.
    \item Temperature, top-$k$ and top-$p$ control diversity.
    \item Guided decoding constrains outputs to valid formats.
    \item Inference optimization is central to deploying LLMs.
\end{itemize}

\end{frame}

%------------------------------------------------
\begin{frame}{Suggested Reading}

\begin{itemize}
    \item Shazeer et al. (2017), \textit{Outrageously Large Neural Networks}
    \item Fedus et al. (2021), \textit{Switch Transformers}
    \item Jiang et al. (2024), \textit{Mixtral of Experts}
    \item Holtzman et al. (2019), \textit{The Curious Case of Neural Text Degeneration}
    \item Leviathan et al. (2023), \textit{Fast Inference from Transformers via Speculative Decoding}
\end{itemize}

\end{frame}

%------------------------------------------------
\begin{frame}
\centering
\Huge Questions?
\end{frame}

\end{document}